\documentclass[11pt]{article}
\usepackage[letterpaper,margin=1in]{geometry}
\usepackage{amsmath,amssymb}
\usepackage[hidelinks]{hyperref}
\setlength{\parindent}{0pt}
\setlength{\parskip}{0.7em}

\title{Literature status: determinants of random unitary pencils}
\author{Run \texttt{fa\_banach\_001} --- lane 12}
\date{11 August 2026}

\begin{document}
\maketitle

\textbf{Status:} \texttt{literature\_already\_answered}.

\section*{Source question}

Jury and Roman's \emph{Determinants of Random Unitary Pencils}
\cite{JR26} asks Conjecture~(1.9) on PDF page~3.  For a tuple
$\mathcal X=(X_1,\ldots,X_g)$ of $k\times k$ matrices, put
\[
 L_{\mathcal X}(\mathcal U)
 =I_k\otimes I_d+\sum_{j=1}^gX_j\otimes U_j,
 \qquad
 \mathbf{rad}(\mathcal X)^2
 =\rho\!\left(\sum_{j=1}^gX_j\otimes\overline{X_j}\right).
\]
The conjecture states that if
$\mathbf{rad}(\mathcal X),\mathbf{rad}(\mathcal Y)<1$, then
\begin{equation}\label{eq:source}
 \lim_{d\to\infty}\int_{U(d)^g}
 \det L_{\mathcal X}(\mathcal U)\,
 \overline{\det L_{\mathcal Y}(\mathcal U)}\,d\mathcal U
 =\det\!\left(
 I_k\otimes I_{k'}-\sum_{j=1}^gX_j\otimes\overline{Y_j}
 \right)^{-1}.
\end{equation}
The same conjecture is restated at the start of Section~3 of the source.

\section*{Identification of the answer}

Jury, van Rensburg, and Roman's later paper \emph{Free versions of the
strong Szeg\H{o} limit theorem} \cite{JvRR26} explicitly identifies and
settles the question.  The sentence immediately preceding Theorem~1.6 on PDF
page~4 says that the theorem ``fully proves the conjecture left open'' in
\cite{JR26}.  The theorem assumes
\[
 \rho\!\left(\sum_{j=1}^g A_j\otimes\overline{A_j}\right)<1,
 \qquad
 \rho\!\left(\sum_{j=1}^g B_j\otimes\overline{B_j}\right)<1,
\]
and concludes precisely the limit in \eqref{eq:source}, with
$(A,B,N,\ell)$ in place of $(\mathcal X,\mathcal Y,d,k')$.  These are exactly
the source's outer-spectral-radius hypotheses after squaring.

The later paper proves more: Theorem~6.2 on PDF page~25 treats quotients of
determinants of monic pencils, and explicitly notes that Theorem~1.6 is the
case $C=D=0$.

\section*{Scope and provenance}

The match is exact, not merely a reformulation inferred by this run.  The
supporting authors include both source authors, cite the source, name its open
conjecture, use the same matrix tuples and spectral-radius domain, and prove
the same determinant limit.  Thus no part of Conjecture~(1.9) remains open in
its stated scope.  This note records literature status only and makes no claim
of new mathematics.

The check used exact-title, arXiv-id, author, ``random unitary pencil,'' and
``outer spectral radius'' searches, followed by direct inspection of the two
papers.  Local copies of both PDFs accompany this note.

\begin{thebibliography}{9}
\bibitem{JR26}
M. T. Jury and G. Roman,
\emph{Determinants of Random Unitary Pencils},
arXiv:2506.04400v1 (2025); J. Funct. Anal. 291 (2026), no.~6, 111555.

\bibitem{JvRR26}
M. T. Jury, L. J. van Rensburg, and G. Roman,
\emph{Free versions of the strong Szeg\H{o} limit theorem},
arXiv:2607.25980v1 (2026).
\end{thebibliography}

\end{document}

